\section{Full Hybrid Sequence}
\label{app:hybrids}

This appendix gives the two simulators of \S\ref{sec:security} in full and expands
\S\ref{sec:hybrids} with the distinguisher reductions, at the deployed operating point
$(w,t,T) = (6,3,18)$ over $d = 128$-bit codes. Throughout, $\mathcal{D}$ is a distinguisher,
$\mathsf{view}_i$ its view in $\mathsf{H}_i$, and
$\varepsilon_i = |\Pr[\mathcal{D}(\mathsf{view}_{i-1}) = 1] - \Pr[\mathcal{D}(\mathsf{view}_i) = 1]|$.

\subsection*{Simulators}

\paragraph{Simulator for a corrupt client.}
$\mathsf{Sim}_{\mathsf{C}}$ receives $\tilde{\mathbf{x}}$ together with
$(\mathsf{out}, \mathcal{L}_{\mathsf{C}})$ and must produce the mOPRF outputs, the
sender-side silent-VOLE transcript, and the polynomials $\{\hat z_i\}_{i \in [N]}$.
\begin{enumerate}\itemsep2pt
\item Sample $r_j \leftarrow \mathcal{D}$ uniformly for $j \in [T]$ as the mOPRF outputs
      and set $\hat X'' = \{ r_j \oplus j \}_{j\in[T]}$.
\item For each $i \in [N]$: if $i \in \mathsf{out}$, sample
      $\hat\ell'_i \leftarrow \mathbb{F}$ and a uniform polynomial $\hat f_i$ of degree
      $\le \theta - 1$ with $\hat f_i(0) = \hat\ell'_i$; otherwise set $\hat f_i = \bot$.
      This $\hat f_i$ is the \emph{label-sharing} polynomial, distinct from the masking
      polynomial $R$ used below --- see the remark at the end of this appendix.
\item Build $\hat z_i$ by fixing its values on the $2T+1$ interpolation points: at points
      indexed by $j \in S_i$ (read from $\mathcal{L}_{\mathsf{C}}$) set
      $\hat z_i(x_j) = \hat f_i(x_j)$; at every other point sample uniformly from
      $\mathbb{F}$; interpolate the unique polynomial of degree $\le 2T$.
\item Emit the sender's correlation messages as uniform vectors of the appropriate length.
\end{enumerate}

\paragraph{Simulator for a corrupt server.}
$\mathsf{Sim}_{\mathsf{S}}$ receives $D$ and the PRF key $k$ and outputs $\bot$, so its
only task is the transcript. The single input-dependent message is
$\Delta U = \Delta u(X)$ with $\Delta u = u_0 - u'_0$. It is worth stating precisely why
this is simulatable, because the natural argument is slightly wrong: $\Delta U$ is
\emph{not} uniform on $\mathbb{F}^{2T+1}$. It lies in the $(T{+}1)$-dimensional image of
the evaluation map on polynomials of degree $\le T$. The simulator therefore samples a
uniform polynomial of degree $\le T$ and outputs its evaluation vector, and simulates the
silent-VOLE receiver messages with the base-VOLE simulator. Since $u'_0$ is uniform,
$u_0 - u'_0$ is uniform \emph{in that subspace}, and the indistinguishability is
statistical rather than computational. This part of the argument does not depend on
$(w,t,T)$ at all --- it depends only on $\deg R = T$, established below to hold by
construction.

\subsection*{Hybrid reductions}

\paragraph{$\mathsf{H}_0 \to \mathsf{H}_1$ (mOPRF).} Replace the masked-OPRF sub-protocol with
$\mathcal{F}^{T}_{\mathsf{mOPRF}}$. Given $\mathcal{D}$ distinguishing these hybrids, build
$\mathcal{B}$ against semi-honest mOPRF security: $\mathcal{B}$ runs the rest of the protocol
honestly, forwards its challenge transcript in place of the sub-protocol messages, and outputs
$\mathcal{D}$'s guess. Hence $\varepsilon_1 \le \varepsilon_{\mathsf{mOPRF}}$.

\paragraph{$\mathsf{H}_1 \to \mathsf{H}_2$ (PRF) --- constant recomputed at the deployed
operating point.} Replace $F_k$ by a uniformly random function. The reduction $\mathcal{B}$
queries its oracle wherever the protocol evaluates $F_k$. The query budget is the subtlety: each
mask restricts the PRF input to a domain of size $2^{w}$, so a distinguisher may enumerate the
entire input table of $T \cdot 2^{w}$ entries, and $\mathcal{B}$ must be allowed that many oracle
calls. At the \emph{deployed} operating point $w = 6$, $T = 18$ this is
\[
T \cdot 2^{w} \;=\; 18 \cdot 2^{6} \;=\; 18 \cdot 64 \;=\; 1152 \;\approx\; 2^{10.2},
\]
giving $\varepsilon_2 \le \varepsilon_{F}(q + 1152)$. This replaces the earlier tuning's
$w = 14$, $T = 64$ figure of $T \cdot 2^{w} = 2^{20}$, which is stale and must not be quoted for
the deployed system: it belongs only to a discussion contrasting against that earlier,
communication-heavier operating point. Omitting the $T \cdot 2^w$ factor at all --- charging the
reduction only $q$ --- is the most common way to understate the bound for this construction,
regardless of which operating point is in force.

\paragraph{$\mathsf{H}_2 \to \mathsf{H}_3$ (VOLE).} Replace the vector ring-OLE sub-protocol with
$\mathcal{F}_{\mathsf{VROLE}}$. Indistinguishability follows from base-VOLE security and the LPN
assumption used by the silent extension~\cite{BoyleEtAl19,RindalSchoppmann21}, so
$\varepsilon_3 \le \varepsilon_{\mathsf{VROLE}} + \varepsilon_{\mathsf{LPN}}$.

\paragraph{$\mathsf{H}_3 \to \mathsf{H}_4$ (statistical masking) --- design constraint verified
satisfied by construction.} For every non-colliding evaluation point $x_j$, replace
$v_1(x_j) = f(x_j) + \bigl(\prod_i (x_j - \yt_i)\bigr) \cdot R(x_j)$
by a uniform element of $\mathbb{F}$.

Fix a record $i$ and let $J = \{j \in [T] : j \notin S_i\}$ be its non-colliding points. At each
such point $\prod_i (x_j - \yt_i) \neq 0$, so the map $R \mapsto (\prod_i(x_j - \yt_i) R(x_j))_{j
\in J}$ is a scaled evaluation map. Since $R$ is uniform of degree $\le T$, its evaluation vector
on any set of at most $T+1$ distinct points is uniform on $\mathbb{F}^{|J|}$ --- the evaluation
map on $\deg R + 1$ points is a bijection onto $\mathbb{F}^{\deg R + 1}$. The masking terms are
therefore jointly uniform, and the substitution is perfect, \emph{provided}
\begin{equation}
\label{eq:h4app}
|J| \;\le\; \deg R + 1 \;=\; T + 1 .
\end{equation}
Here $|J| \le T$ always, so~\eqref{eq:h4app} holds unconditionally given $\deg R = T$. This is a
design constraint, not an accident, and we do not merely assert it: the masking polynomial $R$ is
instantiated in the implementation with $\deg R$ set equal to the sub-sample count, so
$\deg R = T$ by construction, at every operating point --- including the deployed $T = 18$, where
constraint~\eqref{eq:h4app} reads $|J| \le 19$ and holds because $|J| \le T = 18$. We flag one
important scoping note: this coupling is a property of the current implementation binding
$\deg R$ to $T$, not a fact about the abstract protocol independent of that binding. If a future
change were to reduce $\deg R$ \emph{independently} of $T$ --- for instance to save some other
resource without also lowering the sub-sample count --- constraint~\eqref{eq:h4app} would no
longer hold automatically and this hybrid step would need to be re-derived at the new degree. The
communication-driven parameter search that took $T: 64 \to 18$ moved $\deg R$ together with $T$,
so it does not endanger this step; we note the coupling explicitly so that it is checked again if
$\deg R$ and $T$ are ever decoupled in a future revision.

Applying this simultaneously at all $N$ records and all $2T+1$ points, the union bound over the
event that any evaluation point coincides with a root gives
$\varepsilon_4 \le N(2T+1)/|\mathbb{F}|$, which is the statistical term of~\eqref{eq:bound} and
the reason $|\mathbb{F}|$ must be chosen with $N$ in view.

\paragraph{$\mathsf{H}_4 \to \mathsf{H}_5$ (labels).} Replace $\ell'_i = F_k(\ell_i)$ by uniform
field elements. Immediate from $\mathsf{H}_2$, in which $F_k$ is already a random function;
$\varepsilon_5 = 0$.

\paragraph{$\mathsf{H}_5 = \mathsf{Sim}_{\mathsf{C}}$.} In $\mathsf{H}_5$ every message the
client sees is either uniform or determined by $(\xt, \mathsf{out}, \mathcal{L}_{\mathsf{C}})$,
which is exactly the simulator's input, so the two distributions are identical by inspection.

\paragraph{Composition across $N$ records.} All $N$ record instances share one receiver
polynomial $u_0$ and one PRF key $k$, across records and sessions. This is why $\mathsf{H}_2$ and
$\mathsf{H}_4$ are stated over the whole database rather than per record: $\mathsf{H}_2$ replaces
the shared $F_k$ once, so its cost is charged once rather than $N$ times, and $\mathsf{H}_4$ is
applied simultaneously at all $N(2T+1)$ points under a single union bound. A per-record
invocation of a single-instance theorem would not establish this, because correlated reuse across
instances is precisely where such compositions fail.

\paragraph{On the two polynomials.} The construction contains two polynomials that must not be
conflated when reading the two paragraphs above. The label-sharing polynomial $f_i$ (used inside
$\mathsf{Sim}_{\mathsf{C}}$'s step 2, and by the $t$-of-$T$ label-reconstruction protocol) has
degree $\le \theta - 1$; it carries the secret-shared label and its degree is set by the
reconstruction threshold, independent of $T$. The masking polynomial $R$ analysed in the
$\mathsf{H}_3 \to \mathsf{H}_4$ step has degree $T$; it carries no label information and its sole
role is to statistically hide non-colliding evaluation points. The $\deg R = T$ coupling
established above concerns only the second polynomial.

Summing the terms gives~\eqref{eq:bound}.
